\documentclass[conference]{IEEEtran}
\usepackage{amsmath}
\usepackage{booktabs}
\usepackage{graphicx}
\usepackage{hyperref}

\title{CHRONOS: Causal Hypergraph-Regulated Orchestration of Networked Offloading with Spiking Intelligence}

\author{\IEEEauthorblockN{George David Tsitlauri}
\IEEEauthorblockA{\textit{Dept. of Informatics \& Telecommunications}\\
\textit{University of Thessaly, Greece}\\
 gdtsitlauri@gmail.com}}

\begin{document}
\maketitle

\begin{abstract}
CHRONOS is a simulation-first research framework for edge and IoT orchestration under conflicting objectives: model accuracy, latency, communication cost, energy use, and deadline compliance. The repository implements a causal/hypergraph-oriented orchestration stack, baseline comparators, ablation tooling, and reproducibility scripts. In the currently committed three-seed results, CHRONOS attains the highest accuracy-related reward but performs poorly on communication cost and deadline violations, which causes it to underperform simpler baselines on the scalarized combined reward. The main contribution of the current repository state is therefore methodological and diagnostic: it exposes a meaningful trade-off surface rather than supporting a claim of outright scheduling superiority.
\end{abstract}

\section{Introduction}
Autonomous edge orchestration is naturally multi-objective. Policies that increase task accuracy or offloading aggressiveness can easily worsen latency, congestion, and deadline compliance. CHRONOS studies this tension through a synthetic research environment that combines causal state modeling, hypergraph-style coordination, and spiking-policy-inspired decision modules.

\section{Repository Scope}
The repository includes:
\begin{itemize}
\item a configurable edge/IoT simulator,
\item CHRONOS orchestration modules,
\item baseline implementations,
\item ablation and reproducibility scripts,
\item result aggregation under \texttt{results/} and \texttt{outputs/}.
\end{itemize}
The current evidence base is simulation-only. No claim is made about production deployment or hardware validation.

\section{Methodology}
CHRONOS is evaluated against multiple scheduling baselines under a shared synthetic environment. The committed summary aggregates three seeds (42, 43, 44) and reports objective-level rewards as well as a scalarized combined reward. This distinction matters because the qualitative story changes depending on whether one optimizes for task-accuracy gain alone or for overall system utility.

\section{Current Results}
Table~\ref{tab:results} summarizes the main committed findings from \texttt{results/summary.json}.

\begin{table}[h]
\centering
\caption{Committed multi-seed summary (mean $\pm$ std).}
\label{tab:results}
\begin{tabular}{lrrrr}
\toprule
Method & Combined & Acc. reward & Comm. reward & Deadline viol. \\
\midrule
CHRONOS & -47.46 $\pm$ 33.80 & 1.784 $\pm$ 0.760 & -435.92 $\pm$ 308.25 & 0.430 $\pm$ 0.308 \\
FedAvg+Greedy & 0.012 $\pm$ 0.008 & 0.031 $\pm$ 0.020 & 0.00 $\pm$ 0.00 & 0.000 $\pm$ 0.000 \\
DRL-Offload & 0.011 $\pm$ 0.006 & 0.029 $\pm$ 0.015 & 0.00 $\pm$ 0.00 & 0.000 $\pm$ 0.000 \\
GNN-Sched & -0.052 $\pm$ 0.080 & 0.037 $\pm$ 0.036 & -0.59 $\pm$ 0.83 & 0.071 $\pm$ 0.100 \\
\bottomrule
\end{tabular}
\end{table}

The results support three restrained conclusions:
\begin{enumerate}
\item CHRONOS is effective at extracting accuracy-related reward.
\item The current policy is communication-heavy and deadline-expensive.
\item Under the present scalarization, simpler baselines remain stronger on the overall combined objective.
\end{enumerate}

\section{Discussion}
These results do not invalidate the CHRONOS design; they show that the current operating point is not yet well calibrated. In practice, the repository is useful as a research environment for studying objective tension, ablation sensitivity, and reward-design effects. It should not yet be framed as a best-in-class orchestration policy.

\section{Limitations}
The evaluation is fully synthetic, the seed count is still modest, and the current scalarized reward may over- or under-penalize communication-heavy behaviors depending on downstream deployment goals. Real edge traces and hardware-in-the-loop validation remain future work.

\section{Why the Current Release Still Stands Up Well}
The present CHRONOS release is still valuable because it surfaces a real and
important systems trade-off rather than hiding it. The code, baselines,
analysis, and artifacts are already in place; the repository's next step is
calibration, not reconstruction. That makes CHRONOS a credible orchestration
research repo even before the operating point is fully optimized.

\section{Conclusion}
CHRONOS is a credible simulation research repository with real implementation depth. Its current contribution is not a solved orchestration benchmark, but a concrete and reproducible study of the trade-offs that arise when accuracy-seeking coordination policies interact with strict systems constraints.

\bibliographystyle{IEEEtran}
\bibliography{chronos_refs}
\end{document}



